Este manual tiene como propósito ser una guía lo más completa posible del material necesario para participar en el Concurso Internacional Universitario de Programación (ICPC, por sus siglas en inglés) en cuánto a la parte matemática se refiere, es cierto que esta competencia es interdisciplinaria y usualmente se busca que los equipos se conformen por gente con un trasfondo fuerte en algoritmia y además en matemáticas.Este libro se enfocará en los conocimientos necesarios de matemáticas que le conviene saber a un concursante por lo que se recomienda para personas estudiantes de matemáticas que quieran iniciarse en la ICPC o personas participantes de ICPC que busquen expandir sus conocimientos matemáticos. 

Aunque este manual tenga una gran carga de teoría y definiciones, la familiarización con cada uno de los contenidos de este manual radicará no sólo en leer los capítulos pero también en intentar resolver los problemas y ejercicios sugeridos, si bien algunos problemas fueron recabados de jueces en línea, muchos de ellos fueron diseñados con mucho cariño con la única finalidad de servir para el lector de este manual.

Es importante mencionar que si bien este manual está ordenado por capítulos, estos no están ordenados por dificultad, si no por área de las matemáticas a las que pertenecen, por lo que el lector no necesita necesariamente seguir el orden del indice.
% Esta parte modificarla para explciar un poco más la relevancia de cada area en competencias.
El primer capítulo se presentan técnicas de álgebra avanzada usualmente impartidas dentro de los cursos de Algebra Lineal y Algebra Moderna, el segundo capítulo se enfocará en la parte de teoría de números, se cubrirán conocimientos impartidos en Teoría de Números I y II, así como propiedades básicas usualmente vistas en la Teoría Análitica de Números. La sección de Combinatoria ocupará técnicas y propiedades de combinatoria básicas


La segunda parte de este manual esta dedicada a las soluciones de cada una de los problemas presentados en este libro con explicación y código fuente así como algunos hints para poder avanzar en caso de que el lector se encuentre atorado y no quiera leer la solución. Se recomienda al lector realmente intentar resolver los problemas antes de leer las soluciones, pues la parte más valiosa de este manual no es en sí mismo los temas explicados pero la aplicación de cada uno de ellos a los problemas presentados donde las soluciones puede que no requieran temas difíciles pero sí mucho ingenio.

%Escribir una nota personal de por que quise escribir este manual
